\slajdid{09_3-s014}
\begin{frame}[label=09_3-s014]
\textbf{\alert{Optimizacija kašnjenja u}\ {\color{DEplava}lancu}\ \alert{CMOS invertora}}
\par\bigskip
U situaciji kada invertor ima neku fiksnu parazitnu kapacitivnost na izlazu, relativno veliku, postavlja se isto pitanje kako dimenzionisati tranzistore da bi kašnjenje prouzrokovano tom kapacitivnošću bilo što je manje moguće.
\par\bigskip\centering\begin{tikzpicture}[x=.85cm,y=.65cm,font=\sitneoznake,>=Latex]\ctikzset{bipoles/length=.6cm}
\draw(0,-.4)--(1,.2)--(0,.8)--cycle;\draw(1.1,.2)circle(.1);\draw(-.6,.2)node[left]{$v_I(t)$}--(0,.2);\draw(1.2,.2)--(2.8,.2)node[right]{$v_O(t)$};\draw(2.2,.2)to[C,l=$C_L$](2.2,-1.4)node[ground]{};\end{tikzpicture}
\par\bigskip\raggedright
Direktan odgovor je da treba povećati širine kanala oba tranzistora, zadržavajući njihov utvrđen odnos, tako da se povećaju strujni kapaciteti, odnosno smanje njihove otpornosti. Međutim na taj način ćemo značajno \alert{povećati i internu kapacitivnost} samog invertora pa će on \alert{u značajnoj meri opteretiti prethodno kolo i izazvati povećanje ukupnog kašnjenja}.
\end{frame}
